\documentclass[10pt,a4paper]{amsart}
\usepackage{amsmath,amssymb,amsthm,graphicx}
% Retained preamble for the Stacks Project chapter divisors.tex (Divisors).
% The chapter file pinned-original-source.tex carries only \input{preamble} on its first line, so no
% preamble text lives inside the chapter itself.  The Stacks Project uses one single shared preamble for
% every chapter; the snapshot stages it verbatim in sources/stacks-chow/preamble.tex.  The block below is
% copied mechanically, line for line, from that file's lines 14-18 (the Xy-pic setup) and lines 154-217
% (the native theorem declarations together with the mathematical macros).  Nothing was added, removed or
% redefined.  The wrapper class and the amsmath/amssymb/amsthm/graphicx base packages are supplied by the
% assembler, so the class/\externaldocument part of the shared preamble is not repeated here.
% For commutative diagrams we use Xy-pic
\usepackage[all]{xy}

% We use 2cell for 2-commutative diagrams.
\xyoption{2cell}
% Theorem environments.
%
\theoremstyle{plain}
\newtheorem{theorem}[subsection]{Theorem}
\newtheorem{proposition}[subsection]{Proposition}
\newtheorem{lemma}[subsection]{Lemma}

\theoremstyle{definition}
\newtheorem{definition}[subsection]{Definition}
\newtheorem{example}[subsection]{Example}
\newtheorem{exercise}[subsection]{Exercise}
\newtheorem{situation}[subsection]{Situation}

\theoremstyle{remark}
\newtheorem{remark}[subsection]{Remark}
\newtheorem{remarks}[subsection]{Remarks}

\numberwithin{equation}{subsection}

% Macros
%
\def\lim{\mathop{\mathrm{lim}}\nolimits}
\def\colim{\mathop{\mathrm{colim}}\nolimits}
\def\Spec{\mathop{\mathrm{Spec}}}
\def\Hom{\mathop{\mathrm{Hom}}\nolimits}
\def\Ext{\mathop{\mathrm{Ext}}\nolimits}
\def\SheafHom{\mathop{\mathcal{H}\!\mathit{om}}\nolimits}
\def\SheafExt{\mathop{\mathcal{E}\!\mathit{xt}}\nolimits}
\def\Sch{\mathit{Sch}}
\def\Mor{\mathop{\mathrm{Mor}}\nolimits}
\def\Ob{\mathop{\mathrm{Ob}}\nolimits}
\def\Sh{\mathop{\mathit{Sh}}\nolimits}
\def\NL{\mathop{N\!L}\nolimits}
\def\CH{\mathop{\mathrm{CH}}\nolimits}
\def\proetale{{pro\text{-}\acute{e}tale}}
\def\etale{{\acute{e}tale}}
\def\QCoh{\mathit{QCoh}}
\def\Ker{\mathop{\mathrm{Ker}}}
\def\Im{\mathop{\mathrm{Im}}}
\def\Coker{\mathop{\mathrm{Coker}}}
\def\Coim{\mathop{\mathrm{Coim}}}

% Boxtimes
%
\DeclareMathSymbol{\boxtimes}{\mathbin}{AMSa}{"02}

%
% Macros for moduli stacks/spaces
%
\def\QCohstack{\mathcal{QC}\!\mathit{oh}}
\def\Cohstack{\mathcal{C}\!\mathit{oh}}
\def\Spacesstack{\mathcal{S}\!\mathit{paces}}
\def\Quotfunctor{\mathrm{Quot}}
\def\Hilbfunctor{\mathrm{Hilb}}
\def\Curvesstack{\mathcal{C}\!\mathit{urves}}
\def\Polarizedstack{\mathcal{P}\!\mathit{olarized}}
\def\Complexesstack{\mathcal{C}\!\mathit{omplexes}}
% \Pic is the operator that assigns to X its picard group, usage \Pic(X)
% \Picardstack_{X/B} denotes the Picard stack of X over B
% \Picardfunctor_{X/B} denotes the Picard functor of X over B
\def\Pic{\mathop{\mathrm{Pic}}\nolimits}
\def\Picardstack{\mathcal{P}\!\mathit{ic}}
\def\Picardfunctor{\mathrm{Pic}}
\def\Deformationcategory{\mathcal{D}\!\mathit{ef}}

% --- appended 2026-10-12 (Chief): the project own definitions of the reference/slogan/history
% display environments, copied mechanically from the shared Stacks preamble lines 8-12
% (\usepackage{verbatim} plus the three \newenvironment lines, which the chapter preamble keeps
% verbatim from the project preamble). Nothing added, removed or redefined: these environments
% carry no mathematical content and the project itself typesets them as comments. Without these
% three lines a counted statement whose own text uses them fails with Environment slogan undefined.
\usepackage{verbatim}
\newenvironment{reference}{\comment}{\endcomment}
\newenvironment{slogan}{\comment}{\endcomment}
\newenvironment{history}{\comment}{\endcomment}
\usepackage{cleveref}
\crefname{originallemma}{original lemma}{original lemmas}
\Crefname{originallemma}{Original lemma}{Original lemmas}
\makeatletter
\newlabel{lemma-effective-Cartier-codimension-1}{{0BCN}{source}{}{sourcepointer-1}{}}
\newlabel{lemma-effective-Cartier-codimension-1@cref}{{[originallemma][0][]0BCN}{[1][0][]source}{}{}{}}
\makeatother
\crefname{originallemma}{original lemma}{original lemmas}
\Crefname{originallemma}{Original lemma}{Original lemmas}
\makeatletter
\newlabel{lemma-effective-Cartier-divisor-is-a-sum}{{0BCP}{source}{}{sourcepointer-2}{}}
\newlabel{lemma-effective-Cartier-divisor-is-a-sum@cref}{{[originallemma][0][]0BCP}{[1][0][]source}{}{}{}}
\makeatother
\crefname{originallemma}{original lemma}{original lemmas}
\Crefname{originallemma}{Original lemma}{Original lemmas}
\makeatletter
\newlabel{lemma-effective-Cartier-in-points}{{0AG8}{source}{}{sourcepointer-3}{}}
\newlabel{lemma-effective-Cartier-in-points@cref}{{[originallemma][0][]0AG8}{[1][0][]source}{}{}{}}
\makeatother
\crefname{originallemma}{original lemma}{original lemmas}
\Crefname{originallemma}{Original lemma}{Original lemmas}
\makeatletter
\newlabel{morphisms-lemma-flat-base-change-scheme-theoretic-image}{{081I}{source}{}{sourcepointer-4}{}}
\newlabel{morphisms-lemma-flat-base-change-scheme-theoretic-image@cref}{{[originallemma][0][]081I}{[1][0][]source}{}{}{}}
\makeatother
\crefname{originallemma}{original lemma}{original lemmas}
\Crefname{originallemma}{Original lemma}{Original lemmas}
\makeatletter
\newlabel{morphisms-lemma-quasi-compact-scheme-theoretic-image}{{01R8}{source}{}{sourcepointer-5}{}}
\newlabel{morphisms-lemma-quasi-compact-scheme-theoretic-image@cref}{{[originallemma][0][]01R8}{[1][0][]source}{}{}{}}
\makeatother
\crefname{originallemma}{original lemma}{original lemmas}
\Crefname{originallemma}{Original lemma}{Original lemmas}
\makeatletter
\newlabel{properties-lemma-immersion-into-noetherian}{{01OX}{source}{}{sourcepointer-6}{}}
\newlabel{properties-lemma-immersion-into-noetherian@cref}{{[originallemma][0][]01OX}{[1][0][]source}{}{}{}}
\makeatother
\title{An effective Cartier divisor extends across a UFD locus}
\author{The Stacks Project Authors}
\date{stacks-project:divisors.tex:local-snapshot-20261012; GNU FDL 1.2 or later}
\begin{document}
\maketitle
\setcounter{subsection}{1}
\noindent\textit{Editorial note.} All delivered mathematical text is the author's own, copied from the pinned Stacks Project chapter file divisors.tex; no statement or proof step was rewritten, reordered or removed.\par\medskip
\noindent\textit{Editorial note.} Carrier window 3: statement and proof both lie inside the window 6000-8000 granted to this lane by the owner lane lane-s13 on 2026-10-14.\par\medskip
\noindent\textit{Editorial note.} The references of the counted spans are recorded as non-load-bearing navigation pointers; the argument that consumes them stays inside the counted proof span.\par\medskip
\noindent\textit{Editorial note.} Preamble: sources/stacks-divisors/retained-preamble.tex (88 lines, verbatim present).\par\medskip
\noindent\textit{Allowed standard prerequisites (editorial).}\begin{itemize}
\item An immersion into a locally Noetherian scheme is quasi-compact. Properties, Lemma properties-lemma-immersion-into-noetherian, used at proof lines 7103-7105. Makes D -> X quasi-compact.
\item Formation of the scheme theoretic image commutes with opens for quasi-compact maps. Morphisms, Lemma morphisms-lemma-quasi-compact-scheme-theoretic-image, used at proof lines 7105-7107. Reduces to X affine.
\item A divisor is Cartier in a point when its ideal is principal there. Lemma lemma-effective-Cartier-in-points, used at proof lines 7109-7112. Reduces further to the local case.
\item Decomposition of an effective Cartier divisor into integral pieces. Lemma lemma-effective-Cartier-divisor-is-a-sum and Lemma lemma-effective-Cartier-codimension-1, used at proof lines 7115-7122. Produces the height one primes and their prime elements.
\end{itemize}\medskip
\ifdefined\hypertarget\hypertarget{sourcepointer-1}{}\fi
\noindent\textit{Source reference.} Original lemma 0BCN: focused-seeds/AGP/sources/stacks-snapshot-20261011/tags.txt, source lines 14286--14286. Navigation pointer into the Stacks Project; the standard fact it supplies is applied inside the counted proof span and the argument resting on it stays there.\par\medskip
\ifdefined\hypertarget\hypertarget{sourcepointer-2}{}\fi
\noindent\textit{Source reference.} Original lemma 0BCP: focused-seeds/AGP/sources/stacks-snapshot-20261011/tags.txt, source lines 14287--14287. Navigation pointer into the Stacks Project; the standard fact it supplies is applied inside the counted proof span and the argument resting on it stays there.\par\medskip
\ifdefined\hypertarget\hypertarget{sourcepointer-3}{}\fi
\noindent\textit{Source reference.} Original lemma 0AG8: focused-seeds/AGP/sources/stacks-snapshot-20261011/tags.txt, source lines 13186--13186. Navigation pointer into the Stacks Project; the standard fact it supplies is applied inside the counted proof span and the argument resting on it stays there.\par\medskip
\ifdefined\hypertarget\hypertarget{sourcepointer-4}{}\fi
\noindent\textit{Source reference.} Original lemma 081I: focused-seeds/AGP/sources/stacks-snapshot-20261011/tags.txt, source lines 10218--10218. Navigation pointer into the Stacks Project; the standard fact it supplies is applied inside the counted proof span and the argument resting on it stays there.\par\medskip
\ifdefined\hypertarget\hypertarget{sourcepointer-5}{}\fi
\noindent\textit{Source reference.} Original lemma 01R8: focused-seeds/AGP/sources/stacks-snapshot-20261011/tags.txt, source lines 2293--2293. Navigation pointer into the Stacks Project; the standard fact it supplies is applied inside the counted proof span and the argument resting on it stays there.\par\medskip
\ifdefined\hypertarget\hypertarget{sourcepointer-6}{}\fi
\noindent\textit{Source reference.} Original lemma 01OX: focused-seeds/AGP/sources/stacks-snapshot-20261011/tags.txt, source lines 2210--2210. Navigation pointer into the Stacks Project; the standard fact it supplies is applied inside the counted proof span and the argument resting on it stays there.\par\medskip
\begin{lemma}
\label{lemma-closure-effective-cartier-divisor}
Let $X$ be a locally Noetherian scheme. Let $U \subset X$ be
an open and let $D \subset U$ be an effective Cartier divisor.
If $\mathcal{O}_{X, x}$ is a UFD for all $x \in X \setminus U$,
then there exists an effective Cartier divisor $D' \subset X$
with $D = U \cap D'$.
\end{lemma}

\begin{proof}
Let $D' \subset X$ be the scheme theoretic image of the morphism $D \to X$.
Since $X$ is locally Noetherian the morphism $D \to X$ is quasi-compact, see
Properties, Lemma \ref{properties-lemma-immersion-into-noetherian}.
Hence the formation of $D'$ commutes with passing to opens in $X$ by
Morphisms, Lemma \ref{morphisms-lemma-quasi-compact-scheme-theoretic-image}.
Thus we may assume $X = \Spec(A)$ is affine.
Let $I \subset A$ be the ideal corresponding to $D'$.
Let $\mathfrak p \subset A$ be a prime ideal corresponding to a point
of $X \setminus U$.
To finish the proof it is enough to show that $I_\mathfrak p$ is generated
by one element, see Lemma \ref{lemma-effective-Cartier-in-points}.
Thus we may replace $X$ by $\Spec(A_\mathfrak p)$, see
Morphisms, Lemma \ref{morphisms-lemma-flat-base-change-scheme-theoretic-image}.
In other words, we may assume that $X$ is the spectrum of a local
UFD $A$. Then all local rings of $A$ are UFD's. It follows that
$D = \sum a_i D_i$ with $D_i \subset U$ an integral effective Cartier divisor,
see Lemma \ref{lemma-effective-Cartier-divisor-is-a-sum}.
The generic points $\xi_i$ of $D_i$ correspond to prime ideals
$\mathfrak p_i \subset A$ of height $1$, see
Lemma \ref{lemma-effective-Cartier-codimension-1}.
Then $\mathfrak p_i = (f_i)$ for some prime element $f_i \in A$
and we conclude that $D'$ is cut out by $\prod f_i^{a_i}$ as desired.
\end{proof}

\end{document}
